\chapter{Proofs of the finite-width frozen-kernel estimates}
\label{app:finite-width-frozen-proofs}

This appendix proves the finite-width frozen-kernel results stated in
Section~\ref{sec:results-finite-width-frozen}. We use the notation introduced in
Chapter~\ref{cha:methodology}. In particular, \(T_0^{(m)}\) denotes the frozen finite-width
tangent operator at initialization, \(T_\infty\) the infinite-width limiting operator,
\(\mathcal H_K\) the truncated Fourier subspace, \(P_K\) the orthogonal projection onto
\(\mathcal H_K\), and
\[
	B_m^{(K)}=P_KT_0^{(m)}P_K,
	\qquad
	\Lambda_K=P_KT_\infty P_K.
\]
Since the input space is the continuum \(S^1\), there is no sampling randomness in this
appendix. The only randomness comes from the finite number of initialized neurons.

\section{One-neuron tangent feature and kernel}
\label{sec:app-finite-width-one-neuron}

Fix one neuron \(r\), and write
\[
	u_r(\theta):=w_r^\top x(\theta)+b_r.
\]
At initialization, \(b_r=0\), so \(u_r(\theta)=w_r^\top x(\theta)\).

We compute the tangent feature of the \(r\)-th neuron contribution
\[
	a_r\sigma(u_r(\theta))
\]
with respect to the parameter block \((a_r,w_r,b_r)\). First,
\[
	\partial_{a_r}\bigl(a_r\sigma(u_r(\theta))\bigr)
	=
	\sigma(u_r(\theta)).
\]
Next, by the chain rule,
\[
	\nabla_{w_r}\bigl(a_r\sigma(u_r(\theta))\bigr)
	=
	a_r\sigma'(u_r(\theta))\nabla_{w_r}u_r(\theta).
\]
Since
\[
	u_r(\theta)=w_r^\top x(\theta)+b_r,
\]
we have
\[
	\nabla_{w_r}u_r(\theta)=x(\theta).
\]
For ReLU,
\[
	\sigma'(z)=\mathbf 1_{\{z>0\}}
\]
away from \(z=0\), which is sufficient almost surely under the Gaussian initialization.
Hence
\[
	\nabla_{w_r}\bigl(a_r\sigma(u_r(\theta))\bigr)
	=
	a_r\mathbf 1_{\{u_r(\theta)>0\}}x(\theta).
\]
Similarly,
\[
	\partial_{b_r}\bigl(a_r\sigma(u_r(\theta))\bigr)
	=
	a_r\sigma'(u_r(\theta))\partial_{b_r}u_r(\theta)
	=
	a_r\mathbf 1_{\{u_r(\theta)>0\}}.
\]

Thus the one-neuron tangent feature is
\begin{equation}
	\psi_r(\theta)
	=
	\begin{pmatrix}
		\sigma(u_r(\theta))                        \\
		a_r\mathbf 1_{\{u_r(\theta)>0\}}\cos\theta \\
		a_r\mathbf 1_{\{u_r(\theta)>0\}}\sin\theta \\
		a_r\mathbf 1_{\{u_r(\theta)>0\}}
	\end{pmatrix}
	\in\mathbb R^4.
	\label{eq:app-finite-width-tangent-feature}
\end{equation}

The corresponding one-neuron kernel is
\begin{equation}
	\Theta^{(r)}(\theta,\eta)
	=
	\psi_r(\theta)^\top \psi_r(\eta).
	\label{eq:app-finite-width-one-neuron-kernel}
\end{equation}
Substituting \eqref{eq:app-finite-width-tangent-feature} gives
\begin{align}
	\Theta^{(r)}(\theta,\eta)
	 & =
	\sigma(u_r(\theta))\sigma(u_r(\eta)) \notag \\
	 & \quad
	+
	a_r^2
	\mathbf 1_{\{u_r(\theta)>0\}}
	\mathbf 1_{\{u_r(\eta)>0\}}
	\bigl(
	\cos\theta\cos\eta+\sin\theta\sin\eta+1
	\bigr).
	\label{eq:app-finite-width-kernel-expanded}
\end{align}
Using
\[
	\cos\theta\cos\eta+\sin\theta\sin\eta=\cos(\theta-\eta),
\]
we obtain
\begin{equation}
	\Theta^{(r)}(\theta,\eta)
	=
	\sigma(u_r(\theta))\sigma(u_r(\eta))
	+
	a_r^2
	\mathbf 1_{\{u_r(\theta)>0\}}
	\mathbf 1_{\{u_r(\eta)>0\}}
	\bigl(\cos(\theta-\eta)+1\bigr).
	\label{eq:app-finite-width-kernel-final}
\end{equation}

Let \(T^{(r)}\) denote the associated one-neuron integral operator,
\begin{equation}
	(T^{(r)}g)(\theta)
	:=
	\int_0^{2\pi}\Theta^{(r)}(\theta,\eta)g(\eta)\,d\mu(\eta).
	\label{eq:app-finite-width-one-neuron-operator}
\end{equation}
Then, by \eqref{eq:method-finite-width-frozen-average},
\[
	T_0^{(m)}=\frac1m\sum_{r=1}^m T^{(r)}.
\]

\section{Projected one-neuron matrices}
\label{sec:app-finite-width-projected-matrices}

Using the ordered basis
\[
	(\phi_0,\phi_{1,c},\phi_{1,s},\dots,\phi_{K,c},\phi_{K,s})
\]
of \(\mathcal H_K\), the entries of \(B_m^{(K)}\) are
\begin{equation}
	(B_m^{(K)})_{pq}
	=
	\langle \phi_p,T_0^{(m)}\phi_q\rangle_{L^2(\mu)}.
	\label{eq:app-finite-width-B-entry}
\end{equation}
Since \(T_0^{(m)}=\frac1m\sum_{r=1}^m T^{(r)}\),
\[
	(B_m^{(K)})_{pq}
	=
	\frac1m\sum_{r=1}^m
	\langle \phi_p,T^{(r)}\phi_q\rangle.
\]
Define
\begin{equation}
	\xi_{r,pq}:=\langle \phi_p,T^{(r)}\phi_q\rangle.
	\label{eq:app-finite-width-xi}
\end{equation}
Then
\begin{equation}
	(B_m^{(K)})_{pq}
	=
	\frac1m\sum_{r=1}^m \xi_{r,pq}.
	\label{eq:app-finite-width-B-average-entry}
\end{equation}

Equivalently, if \(Z_r\) denotes the \(d\times d\) random matrix defined by
\[
	(Z_r)_{pq}:=\xi_{r,pq},
\]
then
\begin{equation}
	B_m^{(K)}
	=
	\frac1m\sum_{r=1}^m Z_r.
	\label{eq:app-finite-width-B-average}
\end{equation}
Since the neurons are initialized independently and identically, the matrices
\(Z_1,\dots,Z_m\) are i.i.d.

\section{Expectation of the projected operator}
\label{sec:app-finite-width-expectation}

\begin{proof}[Proof of \eqref{eq:results-finite-width-mean}]
	From \eqref{eq:app-finite-width-B-average},
	\[
		\mathbb E[B_m^{(K)}]
		=
		\frac1m\sum_{r=1}^m \mathbb E[Z_r]
		=
		\mathbb E[Z_1].
	\]
	The \((p,q)\)-entry is
	\[
		\mathbb E[(Z_1)_{pq}]
		=
		\mathbb E[\xi_{1,pq}]
		=
		\mathbb E[\langle \phi_p,T^{(1)}\phi_q\rangle].
	\]
	By linearity of expectation,
	\[
		\mathbb E[\xi_{1,pq}]
		=
		\langle \phi_p,\mathbb E[T^{(1)}]\phi_q\rangle
		=
		\langle \phi_p,T_\infty\phi_q\rangle.
	\]
	Since \(T_\infty\) is diagonal in the Fourier basis,
	\[
		\langle \phi_p,T_\infty\phi_q\rangle
		=
		\lambda_q\delta_{pq}.
	\]
	Therefore
	\[
		\mathbb E[B_m^{(K)}]=\Lambda_K.
	\]
\end{proof}

\section{Sub-exponential bound for one projected entry}
\label{sec:app-finite-width-subexp}

We next control \(\xi_{r,pq}\). By definition,
\begin{equation}
	\xi_{r,pq}
	=
	\int_0^{2\pi}\int_0^{2\pi}
	\phi_p(\theta)
	\Theta^{(r)}(\theta,\eta)
	\phi_q(\eta)
	\,d\mu(\eta)\,d\mu(\theta).
	\label{eq:app-finite-width-xi-integral}
\end{equation}
Taking absolute values gives
\[
	|\xi_{r,pq}|
	\le
	\int_0^{2\pi}\int_0^{2\pi}
	|\phi_p(\theta)|
	|\Theta^{(r)}(\theta,\eta)|
	|\phi_q(\eta)|
	\,d\mu(\eta)\,d\mu(\theta).
\]
Since every retained Fourier basis function satisfies
\[
	\|\phi_p\|_\infty\le \sqrt2,
\]
we obtain
\[
	|\xi_{r,pq}|
	\le
	2
	\int_0^{2\pi}\int_0^{2\pi}
	|\Theta^{(r)}(\theta,\eta)|
	\,d\mu(\eta)\,d\mu(\theta).
\]

We now bound the kernel. Since
\[
	u_r(\theta)=w_r^\top x(\theta),
	\qquad
	\|x(\theta)\|=1,
\]
Cauchy--Schwarz gives
\[
	|u_r(\theta)|\le \|w_r\|.
\]
Therefore
\[
	0\le \sigma(u_r(\theta))\le |u_r(\theta)|\le \|w_r\|,
\]
and hence
\[
	0\le \sigma(u_r(\theta))\sigma(u_r(\eta))\le \|w_r\|^2.
\]
For the derivative and bias term, we use
\[
	\mathbf 1_{\{u_r(\theta)>0\}}\mathbf 1_{\{u_r(\eta)>0\}}\le 1
\]
and
\[
	0\le \cos(\theta-\eta)+1\le 2.
\]
Thus
\[
	a_r^2
	\mathbf 1_{\{u_r(\theta)>0\}}
	\mathbf 1_{\{u_r(\eta)>0\}}
	(\cos(\theta-\eta)+1)
	\le
	2a_r^2.
\]
Using \eqref{eq:app-finite-width-kernel-final}, we conclude that
\[
	|\Theta^{(r)}(\theta,\eta)|
	\le
	\|w_r\|^2+2a_r^2.
\]
Since \(\mu\) is a probability measure,
\[
	\int\!\!\int
	(\|w_r\|^2+2a_r^2)\,d\mu(\eta)d\mu(\theta)
	=
	\|w_r\|^2+2a_r^2.
\]
Therefore
\begin{equation}
	|\xi_{r,pq}|
	\le
	2\|w_r\|^2+4a_r^2.
	\label{eq:app-finite-width-xi-bound}
\end{equation}

Define
\[
	Y_r:=2\|w_r\|^2+4a_r^2.
\]
Then
\[
	|\xi_{r,pq}|\le Y_r
	\qquad\text{for all }p,q.
\]

We now show that \(Y_r\) is sub-exponential with an explicit parameter. Since
\(w_r\sim \mathcal N(0,I_2)\), its squared Euclidean norm is the sum of two independent
standard Gaussian squares, and therefore \(\|w_r\|^2\) follows a chi-squared distribution
with two degrees of freedom \(\|w_r\|^2\sim \chi^2_2\).

Similarly, since \(a_r\sim \mathcal N(0,1)\), its square follows a chi-squared distribution
with one degree of freedom \(a_r^2\sim \chi^2_1\) and these variables are independent, for \(t<1/8\),
\[
	\mathbb E[e^{tY_r}]
	=
	\mathbb E[e^{2t\|w_r\|^2}]
	\mathbb E[e^{4ta_r^2}].
\]
If \(X\sim\chi^2_\nu\), then
\[
	\mathbb E[e^{sX}]=(1-2s)^{-\nu/2},
	\qquad s<1/2.
\]
Hence
\[
	\mathbb E[e^{2t\|w_r\|^2}]=(1-4t)^{-1},
	\qquad
	\mathbb E[e^{4ta_r^2}]=(1-8t)^{-1/2}.
\]
Thus
\begin{equation}
	\mathbb E[e^{tY_r}]
	=
	(1-4t)^{-1}(1-8t)^{-1/2},
	\qquad t<1/8.
	\label{eq:app-finite-width-Y-mgf}
\end{equation}

Choosing \(t=1/16\), we obtain
\[
	\mathbb E[e^{Y_r/16}]
	=
	\left(\frac34\right)^{-1}
	\left(\frac12\right)^{-1/2}
	=
	\frac43\sqrt2
	<2.
\]
Therefore, by the standard definition of the \(\psi_1\)-norm,
\begin{equation}
	\|Y_r\|_{\psi_1}\le 16.
	\label{eq:app-finite-width-Y-psi1}
\end{equation}
Since \(|\xi_{r,pq}|\le Y_r\),
\[
	\|\xi_{r,pq}\|_{\psi_1}\le 16.
\]
Finally,
\[
	\|\xi_{r,pq}-\mathbb E[\xi_{r,pq}]\|_{\psi_1}
	\le
	\|\xi_{r,pq}\|_{\psi_1}+|\mathbb E[\xi_{r,pq}]|.
\]
Using \(|\mathbb E[\xi_{r,pq}]|\le \mathbb E|\xi_{r,pq}|\) and
\(\mathbb E|X|\le \|X\|_{\psi_1}\), we obtain
\begin{equation}
	\|\xi_{r,pq}-\mathbb E[\xi_{r,pq}]\|_{\psi_1}\le 32.
	\label{eq:app-finite-width-centered-psi1}
\end{equation}

\section{Entrywise and blockwise concentration}
\label{sec:app-finite-width-concentration}

\begin{proof}[Proof of \eqref{eq:results-finite-width-max-bound}]
	Fix \(p,q\), and define
	\[
		X_{r,pq}:=\xi_{r,pq}-\mathbb E[\xi_{r,pq}].
	\]
	Then \(X_{1,pq},\dots,X_{m,pq}\) are i.i.d., centered, and by
	\eqref{eq:app-finite-width-centered-psi1},
	\[
		\|X_{r,pq}\|_{\psi_1}\le 32.
	\]
	Moreover,
	\[
		(B_m^{(K)})_{pq}-(\Lambda_K)_{pq}
		=
		\frac1m\sum_{r=1}^m X_{r,pq}.
	\]

	We apply the standard Bernstein inequality for centered i.i.d. sub-exponential random
	variables: if \(X_1,\dots,X_m\) satisfy \(\|X_i\|_{\psi_1}\le L\), then there exists a
	universal constant \(c>0\) such that
	\[
		\mathbb P\!\left(
		\left|\frac1m\sum_{r=1}^m X_r\right|\ge \varepsilon
		\right)
		\le
		2\exp\!\left(
		-cm\min\!\left(\frac{\varepsilon^2}{L^2},\frac{\varepsilon}{L}\right)
		\right).
	\]
	Applying this with \(L=32\), we obtain
	\begin{equation}
		\mathbb P\!\left(
		\left|(B_m^{(K)})_{pq}-(\Lambda_K)_{pq}\right|
		\ge \varepsilon
		\right)
		\le
		2\exp\!\left(
		-cm\min\!\left(\frac{\varepsilon^2}{32^2},\frac{\varepsilon}{32}\right)
		\right).
		\label{eq:app-finite-width-entry-concentration}
	\end{equation}

	There are \(d^2=(2K+1)^2\) entries. By the union bound,
	\[
		\mathbb P\!\left(
		\|B_m^{(K)}-\Lambda_K\|_{\max}\ge \varepsilon
		\right)
		\le
		\sum_{p,q=1}^d
		\mathbb P\!\left(
		\left|(B_m^{(K)})_{pq}-(\Lambda_K)_{pq}\right|
		\ge \varepsilon
		\right).
	\]
	Using \eqref{eq:app-finite-width-entry-concentration},
	\[
		\mathbb P\!\left(
		\|B_m^{(K)}-\Lambda_K\|_{\max}\ge \varepsilon
		\right)
		\le
		2d^2
		\exp\!\left(
		-cm\min\!\left(\frac{\varepsilon^2}{32^2},\frac{\varepsilon}{32}\right)
		\right).
	\]
	Since \(d=2K+1\), this proves \eqref{eq:results-finite-width-max-bound}.
\end{proof}

% \section{Operator-norm bound}
% \label{sec:app-finite-width-op-proof}
%
% \begin{proof}[Proof of \eqref{eq:results-finite-width-op-bound}]
% 	For every \(d\times d\) matrix \(M\),
% 	\[
% 		\|M\|_{\mathrm{op}}\le d\|M\|_{\max}.
% 	\]
% 	Therefore, if
% 	\[
% 		\|B_m^{(K)}-\Lambda_K\|_{\mathrm{op}}\ge \delta,
% 	\]
% 	then
% 	\[
% 		\|B_m^{(K)}-\Lambda_K\|_{\max}\ge \frac{\delta}{d}.
% 	\]
% 	Hence
% 	\[
% 		\mathbb P\!\left(
% 		\|B_m^{(K)}-\Lambda_K\|_{\mathrm{op}}\ge \delta
% 		\right)
% 		\le
% 		\mathbb P\!\left(
% 		\|B_m^{(K)}-\Lambda_K\|_{\max}\ge \frac{\delta}{d}
% 		\right).
% 	\]
% 	Applying \eqref{eq:results-finite-width-max-bound} with \(\varepsilon=\delta/d\) yields
% 	\[
% 		\mathbb P\!\left(
% 		\|B_m^{(K)}-\Lambda_K\|_{\mathrm{op}}\ge \delta
% 		\right)
% 		\le
% 		2d^2
% 		\exp\!\left(
% 		-cm\min\!\left(
% 			\frac{\delta^2}{32^2d^2},
% 			\frac{\delta}{32d}
% 			\right)
% 		\right).
% 	\]
% 	which proves \eqref{eq:results-finite-width-op-bound}.
% \end{proof}
% \chapter{\label{cha:glossary}Glossary}
%
% In this appendix we give an overview of frequently used terms and
% abbreviations.
%
% \begin{description}
% 	\item[foo:] \ldots
% 	\item [bar:] \ldots
% \end{description}
%
